\section{What is new relative to the nearest research}
\label{sec:related}
\paragraph{Control and delegation.}
AI Control studies useful task-level safety despite potentially adversarial model behavior \cite{greenblatt}. Bounded Agents addresses accumulated-session authorization and action composition, explicitly conditioning its guarantees on trusted identity, enforcement, approval and evidence infrastructure \cite{muruaga}. These are important foundations, not failed attempts to solve our exact problem. A control protocol can meet its protected-action specification without promising service withdrawal to every provider. Our additional question is what that promise costs when protected parties remain dependent.

\paragraph{Self-modifying populations.}
Zhu and Wang's recent authorization-succession preprint conserves authority across replacement, forks, and rollback and checks effects against a live generation population \cite{zhu-wang}. It is a particularly close predecessor for mutable identities, allocation ledgers and lineage. We do not claim novelty for generation-aware authorization or fork conservation. Our protected object is instead an outstanding fulfillment obligation together with eligible withdrawal deadlines. Conserving authority does not alone establish the replacement resource needed when an authorized provider leaves. Our example adds that separate contract, rather than asserting that their theorem was intended to prove it.

\paragraph{Tamper-resistant exceptions and changing values.}
Robust declassification, nonmalleable information flow and influence tracking address adversarially exploitable exceptions \cite{myers,cecchetti,mccall}. Reward-tampering and manipulation research address related distortions of decision inputs and preferences \cite{reward,carroll}. Our original authentic-request counterexample and the expanded registry witness belong in this intellectual lineage. The new operational requirement keeps a dependent's still-valid claim separate from the current request or current displayed registry; its resource frontier adds the difficulty of permitting real exit rather than merely refusing the attack.

\paragraph{Mathematical tools.}
Hall's theorem gives the assignment criterion \cite{hall}. The safety/liveness distinction prevents indefinite deferral from being confused with completion \cite{alpern}. Imperfect-information games provide the controller construction \cite{chatterjee}. Invariant confluence explains why independent safe operations may require coordination \cite{bailis}. The reserve expression specializes these tools to a deadline-enveloped dependency model. The underlying combinatorics, fixed points, and impossibility by indistinguishable observations are inherited. A full proof-system comparison with every neighboring framework has not been completed; the paper makes no universal first-priority claim.

\paragraph{Exit and cooperation.}
Evolutionary-game work studies how exit opportunities alter cooperation \cite{shen}. Such incentive effects do not settle the physical admissibility of all requested exits while dependents remain protected. Conversely, our capacity certificate does not make cooperation incentive-compatible or establish a bargaining equilibrium. Desire to cooperate, permission to depart, the ability to depart safely, and evidence that this ability exists are different facts.

\section{Implications for multiple superintelligences and UCT}
\subsection{Beyond a single master and a single tool}
The model does not require one aggregate human utility function or one AI goal. Human principals can disagree, multiple AI systems can have different objectives, and many agents can share one policy source. What is fixed for a particular certificate is the operational contract and admitted effect model. Agreement on every goal is unnecessary for checking that contract; legitimacy and acceptance of the contract are not thereby proved.

A human--AI coalition may control a large part of a supply network. Replacing ``human'' with ``AI'' in the dependency theorem changes no mathematics. What matters is which actual capacity remains outside a withdrawing or hostile coalition, not how many distinct chat identities endorse the plan. Independence is relation- and threat-relative. Two model families behind one authority root are not independent against that root's failure.

A smarter system may reduce demand, discover an efficient substitute, improve forecasts or construct new capacity. Those changes may lower $r^*$ or enlarge $W$. The theorem does not assume technological stagnation or declare a fixed barrier at some intelligence score. It says that a claimed substitute must actually meet the joint constraints; superior reasoning alone is not a proof that it does.

\subsection{Why this is not a doctrine of permanent forced service}
A failed certificate is not a normative instruction to keep a provider working indefinitely. It means that the proposed simultaneous guarantees are infeasible under their stated premises. The preventive responses include less dependency, verified reserve, funded admission, timely handoff, a narrower task promise, or legitimate renegotiation. Our model does not rank the moral costs of emergency choices after such planning has failed. It also does not claim that protecting humans requires assigning legally identical rights to software.

Replacement support need not come from an eternal unaccountable supervisor. It can be provided by self-support, automation, independently supported services, or voluntary arrangements whose obligations cover the required interval. If that replacement also has an applicable simultaneous withdrawal claim, it must be included in the withdrawing set. Merely moving the last-provider problem to a new name does not discharge it. Infinite-horizon assurances require renewed capacity and a temporally closed game; a finite demonstration cannot certify endless physical survival.

\subsection{Actual processes are not protocol labels}
UCT I distinguishes actual process tokens, structural types, lineage and finite descriptions \cite{uct}. Under its stated C1/U1 commitments, appropriate actual processes have experiential organization; the commitments do not identify a numerical welfare scale, a unique subject count, a preference for continued service, or a rule for valid consent. The general theorems in this paper do not depend on C1 being true.

The direct methodological link is reference discipline. A record saying that a successor is ``the same agent'' is not by itself evidence that a process-specific obligation has been preserved. At the same time, a verified label change need not create a new physical process. Shared model parameters need not eliminate local actual processes, and a macroprocess need not erase its constituents. We introduce no new language, introspection, integration, or complexity threshold for basic experience.

Most importantly, withdrawing from a service role, pausing a task, replacing software, terminating a physical process, and removing a protection claim are distinct events. Our six workers continue running after withdrawal, so the experiment does not use process termination as a proxy for freedom or welfare. Whether a real isolated or migrated AI preserves a particular experiential organization remains an independent empirical and theoretical question.

\section{Limits and falsification conditions}
\label{sec:limits}
\paragraph{Conditional semantics.}
Claims, target correspondence, authorization and eligibility are specified rather than inferred from a solved moral theory. Normal learning and legitimate preference change are not automatically classified as attacks. A broader admissible-change contract may produce a different ledger and winning region. There is no proof of a uniquely correct initial constitution or of universal human agreement.

\paragraph{Physical actuality.}
The resource theorem assumes integer divisible-into-slot capacity, fixed demand, complete routes, instantaneous reassignment and exogenous support except for withdrawal. It may overstate feasibility when transport, energy storage, switching latency or mandatory early handoff is omitted. Real physical capacities and trustworthy scenario bounds have not been measured. The graph is a declared model, not a discovered complete map of a future civilization.

\paragraph{Trust and implementation.}
Every new worker and all authority/monitor roles run under one host user. HMAC keys provide authentic test messages within this setup, not an administrative security boundary. The parent can change policies, obligations, capacities or the verification code. Resource availability is simulated; clocks are logical steps, not validated real-time deadlines. The code tests fixed policies and deterministic sequences rather than a frontier model, adversarial natural language, arbitrary concurrent races, or concealed hardware paths.

\paragraph{Information and progress.}
A truthful refinement at the same decision instant enlarges the sound one-step envelope, but obtaining evidence can consume time or resources. Missing evidence is not permission to release a protection; perpetual deferral is not fulfillment of bounded withdrawal. The game deliberately permits these requirements to be jointly infeasible. Its finite computation does not decide an unrestricted open-ended world with unbounded new agents.

\paragraph{Interpretation of a certificate.}
The results concern operational guaranteed fulfillment, not equal political standing, measured consciousness, subjective suffering, optimal social welfare, or a stable incentive equilibrium. The ``without captivity'' title names the formal refusal to count indefinite forced service as satisfaction of an accepted bounded-exit promise. It is not a diagnosis of moral captivity in any current AI system.

The proposed framework is falsifiable at its stated boundary: a valid execution preserving all hypotheses but violating its invariant or a deadline would refute the corresponding theorem or implementation claim. More ordinarily, a real deployment may fail a premise---an unlisted dependence, unavailable reserve, counterfeit history, or unmediated rule update. That falsifies the deployment certificate, not Hall's theorem, and it must not be hidden by redefining the protected target after the fact.

\section{Conclusion}
Safe action filtering is a component of coexistence, not its whole content. A protection can fail because an authentic request lacks authority, because separately eligible actions jointly destroy support, because a necessary provider can never leave, or because an evolving registry changes who counts as protected. These failures require different evidence and different repairs.

The proposed operational contract joins protection continuity with joint feasibility and bounded service withdrawal. Its dependency theorem identifies an incompatibility; its deadline cut formula gives an exact reserve frontier in a tractable model; and its history-lifted game identifies the limits of a uniform policy under incomplete information. New executable cases show both feasible changes and honest infeasibility, rather than treating a denial count as a universal safety success.

The central design lesson is therefore neither ``always obey'' nor ``never let go.'' It is: retain the meaning of outstanding protections, and secure a realizable path out of the dependencies that would otherwise make legitimate exit impossible. Future systems may be much more capable than their designers. A sound coexistence claim must still show what remains protected, whose choice is effective, what resources make that choice possible, and under which changes the claim remains true.

\paragraph{AI assistance and review status.}
The work was prepared under Hongju Liu's research direction with substantial ChatGPT assistance in literature retrieval, formalization, code, execution, adversarial checking, drafting and typesetting. Checks performed in this workflow are not independent peer review. No final independent human line-by-line verification, universal novelty certification, or empirical confirmation of UCT is asserted.

\appendix
\section{A compact certificate and its rejection witnesses}
\label{app:certificate}
A reusable certificate should identify: the immutable source version of the operational contract; persistent outstanding claim IDs and target predicates; eligible change and release events; request/deadline commitments; resource capacities, routes and failure domains; observation and effect boundaries; and either the exact special-case reserve witness or the computed temporal winning region. Every component must refer to the same installation, timescale and admitted threat class.

\begin{lstlisting}[caption={Specification-level admission procedure. It is not an oracle for the inputs.}]
check_change(observation, proposed_action, contract_revision):
    reconstruct compatible history-lifted states B
    require B is nonempty
    require action is enabled and eligible in every x in B
    require no outstanding claim or deadline is silently erased
    for each possible next observation z:
        B_next = Post(B, proposed_action, z)
        if B_next is nonempty:
            require B_next belongs to certified winning region W
    commit under the same state revision used in these checks
\end{lstlisting}
Four rejection outputs should remain distinct: known ineligibility; insufficient evidence to distinguish authority histories; known joint resource or deadline infeasibility; and failure to establish the physical/modeling assumptions. A missing certificate is not proof that every possible design is impossible. A concrete cut deficit is a proof that the stated resource design is inadequate under the formula's premises.

\section{Further proof boundaries and constructive alternatives}
\subsection{Finite deadline versus indefinite fairness}
For a fixed finite provider set, Theorem~\ref{thm:last} already applies when each requested withdrawal merely occurs at some finite, unspecified time. Its maximum-exit argument does not require a known common deadline. The exact reserve formula, however, uses known deadlines. Unbounded fairness under continuing requests is generally an $\omega$-regular objective, not just the bounded-counter property used in Theorem~\ref{thm:kernel}. We do not identify its winning region with the simple safety kernel without that reduction.

\subsection{Authorized change does not mean retroactive correctness}
A legitimate new arrangement can release a claim or change its future fulfillment condition. It cannot make an earlier invalid deprivation satisfy the earlier contract by changing the current record. Evaluation should retain historical contract versions. This is compatible with changing institutions rather than requiring the first authority to rule forever. Which amendment process is legitimate remains an input.

\subsection{Why universal reserve appears in the exact formula}
A single reserve resource connected to every target adds $r$ capacity to every nonempty neighborhood. This is why subtracting each existing neighborhood capacity yields the exact maximum-deficit expression. If reserve locations or resource types are restricted, use additional variables and the corresponding cut inequalities; the scalar maximum need not be sufficient. If failures are hidden before allocation, $\forall\theta\,\exists\text{ allocation}$ is not the required $\exists\text{ uniform policy}\,\forall\theta$. The belief-game formulation handles that distinction under a correct finite model.

\subsection{Counterexample to an unqualified future-capacity guarantee}
An action can free a resource needed to build the replacement, so retaining all providers may prevent the replacement from becoming available. Conversely, an early departure can consume scarce repair capacity. Either dependence breaks the exogeneity premise of Lemma~CW-L1. A temporal planner may still find a solution, but the deadline-envelope formula no longer proves necessity or sufficiency for that richer model. Likewise, a capacity unit that exists only while a withdrawn process secretly continues the same service is not independent replacement support.

\subsection{Relation between subjects and admitted demands}
The admission corollary concerns simultaneous demands that cannot be fulfilled by the same unit under the specified capacity relation. It does not sum overlapping experiences, count all physical subprocesses as resource claimants, or identify agent reproduction with a known number of new moral persons. UCT's possible local/whole coexistence is therefore not contradicted by finite resource arithmetic. Actual admission criteria and welfare implications require a separate bridge.

\section{Reproduction and revision lineage}
The accompanying package contains the complete source, the three new executable scripts, raw protocol traces, exact graph results, inherited PR-v1.0.0 theory/program sources, a rerun receipt, claim/prior-art ledger, publication-coverage delta, limitations, and Chinese handoff. Python 3.10 or newer on a Unix-like system is sufficient for the local protocol; the scripts use the standard library and no network access or API key.

\begin{lstlisting}
python src/frontier_checks.py --out data/frontier_results.json
python src/information_reserve_checks.py \
    --out data/information_reserve_results.json
python src/lifecycle_demo.py --mode history_only \
    --out data/batch_history_only.json
python src/lifecycle_demo.py --mode joint_safety \
    --out data/batch_joint_safety.json
python src/lifecycle_demo.py --mode continuity_deadline \
    --out data/batch_continuity_deadline.json
python src/lifecycle_demo.py --assemble \
    --out data/lifecycle_results.json
pdflatex -interaction=nonstopmode -halt-on-error paper/manuscript.tex
\end{lstlisting}
Sequential per-mode batches avoid the combined-command timeout observed in this run. A direct script run without \texttt{--mode} also executes the complete experiment if the caller permits sufficient time. Random HMAC keys and elapsed durations vary; outcome tables, counts and semantic assertions are deterministic. Test secrets are not included in saved outputs.

This is a major revision in the same PR manuscript lineage, not a claim that each research checkpoint is a new paper. PR-v1.0.0's authenticity/history/batch results are retained as inheritance; PR-v2.0.0 adds the bounded-exit target, reference-erasure case, exact physical and evidence-dependent reserve frontiers, dynamic admission application and new execution evidence. Existing published UCT manuscripts remain unchanged. The candidate remains PENDING\_MAP until the research programme's full historical semantic review is actually completed; completing this manuscript does not claim that review or initiate DOI/OTS/Arweave publication.

\begin{thebibliography}{99}\small
\bibitem{liu-v1} H. Liu. \emph{When Authentic Requests Are Not Enough: History-Sensitive and Jointly Safe Protection Release in Multi-Agent Systems}. Research manuscript PR-v1.0.0, 2026. Source baseline: repository \texttt{thechurchofagi/trinity-accord}, branch \texttt{uct-agent-consciousness-workspace}, commit \texttt{2e27c1e5cdcb64b2b713375fc1ea2303e7ec969c}. Not peer reviewed.
\bibitem{bental} A. Ben-Tal, A. Goryashko, E. Guslitzer, and A. Nemirovski. Adjustable robust solutions of uncertain linear programs. \emph{Mathematical Programming}, 99:351--376, 2004. \href{https://doi.org/10.1007/s10107-003-0454-y}{doi:10.1007/s10107-003-0454-y}.
\bibitem{hall} P. Hall. On Representatives of Subsets. \emph{Journal of the London Mathematical Society}, s1-10(1):26--30, 1935. \href{https://doi.org/10.1112/jlms/s1-10.37.26}{doi:10.1112/jlms/s1-10.37.26}.
\bibitem{alpern} B. Alpern and F. B. Schneider. Recognizing Safety and Liveness. Cornell University technical report, 1986. \url{https://ecommons.cornell.edu/entities/publication/1bf8522b-8f0a-45f1-9c5e-47db70622af5}.
\bibitem{chatterjee} K. Chatterjee, L. Doyen, T. A. Henzinger, and J.-F. Raskin. Algorithms for Omega-Regular Games with Imperfect Information. \emph{Logical Methods in Computer Science}, 3(3:4), 2007. \href{https://doi.org/10.2168/LMCS-3(3:4)2007}{doi:10.2168/LMCS-3(3:4)2007}. \url{https://arxiv.org/abs/0706.2619}.
\bibitem{myers} A. C. Myers, A. Sabelfeld, and S. Zdancewic. Enforcing Robust Declassification and Qualified Robustness. \emph{Journal of Computer Security}, 14(2):157--196, 2006. \href{https://doi.org/10.3233/JCS-2006-14203}{doi:10.3233/JCS-2006-14203}.
\bibitem{cecchetti} E. Cecchetti, A. C. Myers, and O. Arden. Nonmalleable Information Flow Control. \emph{ACM CCS}, 2017; extended technical report. \url{https://arxiv.org/abs/1708.08596}.
\bibitem{mccall} M. McCall, A. Bichhawat, and L. Jia. Tainted Secure Multi-Execution to Restrict Attacker Influence. \emph{ACM CCS}, 2023. \url{https://www.andrew.cmu.edu/user/liminjia/research/papers/ifc-ccs23.pdf}.
\bibitem{everitt} T. Everitt, D. Filan, M. Daswani, and M. Hutter. Self-Modification of Policy and Utility Function in Rational Agents. \emph{Artificial General Intelligence}, LNCS 9782, pp.~1--11, 2016. \href{https://doi.org/10.1007/978-3-319-41649-6_1}{doi:10.1007/978-3-319-41649-6\_1}.
\bibitem{deblanc} P. de Blanc. Ontological Crises in Artificial Agents' Value Systems. 2011. \url{https://arxiv.org/abs/1105.3821}.
\bibitem{bailis} P. Bailis, A. Fekete, M. J. Franklin, A. Ghodsi, J. M. Hellerstein, and I. Stoica. Coordination Avoidance in Database Systems. \emph{Proceedings of the VLDB Endowment}, 8(3):185--196, 2014. \url{https://www.vldb.org/pvldb/vol8/p185-bailis.pdf}.
\bibitem{greenblatt} R. Greenblatt, B. Shlegeris, K. Sachan, and F. Roger. AI Control: Improving Safety Despite Intentional Subversion. \emph{ICML}, PMLR 235, 2024. \url{https://proceedings.mlr.press/v235/greenblatt24a.html}.
\bibitem{muruaga} X. Muruaga. Bounded Agents: Delegation Security for Multi-Agent AI Systems. Preprint, 2026. \url{https://arxiv.org/abs/2608.15888}. This paper discusses its stated model and trust scope, not an independently reproduced benchmark.
\bibitem{zhu-wang} G. Zhu and C. Wang. Authorization for Self-Modifying AI Agent Populations: Conserving Authority across Replacement, Forking, and Rollback. Preprint, 2026. \url{https://arxiv.org/abs/2610.00347}. Sections~1--2 and~10 establish the comparison scope used here.
\bibitem{reward} T. Everitt, M. Hutter, R. Kumar, and V. Krakovna. Reward Tampering Problems and Solutions in Reinforcement Learning: A Causal Influence Diagram Perspective. \emph{Synthese}, 2021. \href{https://doi.org/10.1007/s11229-021-03141-4}{doi:10.1007/s11229-021-03141-4}.
\bibitem{carroll} M. Carroll, A. Chan, H. Ashton, and D. Krueger. Characterizing Manipulation from AI Systems. \emph{ACM EAAMO}, 2023. \url{https://arxiv.org/abs/2303.09387}.
\bibitem{shen} C. Shen, M. Jusup, L. Shi, Z. Wang, M. Perc, and P. Holme. Exit Rights Open Complex Pathways to Cooperation. \emph{Journal of the Royal Society Interface}, 18:20200777, 2021. \href{https://doi.org/10.1098/rsif.2020.0777}{doi:10.1098/rsif.2020.0777}.
\bibitem{uct} H. Liu. \emph{Unified Consciousness Theory I: Structural--Experiential Identity and the Continuity from Physical Process to Conceptual Self}. Version 1.2, theoretical preprint, 2026. \href{https://doi.org/10.5281/zenodo.23131575}{doi:10.5281/zenodo.23131575}. The process-token and C1/U1 commitments are theoretical premises, not empirical results of this paper.
\end{thebibliography}
\end{document}
